% a06_BALtorus_d3: part of Step 3 of the proof of thm:a06:T1 (input from sec/a06_BALtorus_c.tex).
% Group (D), second half: the bounds (VAC-flat)_k and (DOM)_k, the collar clause of Ba2 (2.57), and the bad localisations of Ba3 p. 422.

\subsubsection*{Step 3(D$_2$): the counterterm moves of \cite[pp.~411--422]{Ba3}, bounds}

\paragraph{Inputs.} (I1) Proposition~1 in the scope Ba{\l}aban states for it: (I1a) the expressions of the expansion \cite[pp.~414--415]{Ba3} (vertices \cite[(1.6)--(1.15)]{Ba3} and products \cite[(1.18)]{Ba3}; propagators $G_k(\Omega,\tilde B)$, $G_k(\Omega_2,\tilde B)$ and the vector $G_k$ on any admissible $\Omega\supset\Omega_2$, the torus included; vacuum expressions included, since vacuum contour graphs are among ``an arbitrary expression in the expansion'', \cite[p.~414]{Ba3}, with $D(G)\ge2$); (I1b) $\delta G_k(\Omega,\Omega_2,\tilde B)$ and \cite[(1.16)]{Ba3} treated as external fields \cite[pp.~414, 419]{Ba3}; (I1c) external fields that are functions of many variables, in particular ``expressions obtained by an integration and a cumulant expansion in $k$th step, where some external legs are contracted by propagators $C^{(k)}_{\Lambda^{(k)}_5}$'' \cite[p.~419]{Ba3}; (I1d) the localisation of the vector legs by smooth partitions of unity \cite[p.~420]{Ba3}, i.e.\ smooth multipliers on single legs from a fixed $C^2$-bounded class, with the $O(1)$ uniform in them, in the domains and in $\tilde B$ \cite[p.~421]{Ba3}. (I2) Norms of the external fields: (I2a) $\delta G$ at depth, from \cite[Thm., (1.10)--(1.12)]{Ba4} in its restricted form made two-sided (\cref{lem:a06A:twosided}): for cubes $\square,\square'$ at depths $\rho_\square,\rho_{\square'}\ge R_0$ from $\Omega'^c$, $\norm{h_\square\delta G_k(\Omega,\Omega')h_{\square'}}_{1,\alpha_0}\le Ce^{-(\delta_0/4)(\rho_\square+\rho_{\square'}+\dist(\square,\square'))}$; (I2b) block-level covariance differences ($C_\Lambda$ against $C^{(k)}$, \cref{lem:a06A:CLdiff}; domain differences of $\Delta^{(k)}$ and $C_\Lambda$, \cref{lem:a06A:domdiff}), lifted by $a_kG_kQ_k^*$ whose $C^{1,\alpha_0}$ regularity is \cite[Thm., (1.9)--(1.10)]{Ba4} and \cite[Prop.~2.2, (2.58)]{Ba2}, vector versions at $B=0$ componentwise; (I2c) the lifted index-$k$ kernels $\beta_kG_kR_k^*C^{(k)}R_kG_k\beta_k$, with the same inputs and the decay of $C^{(k)}$ (\cref{lem:a06A:blockinv}(c)). (I3) The size of the finite mass vertices: in scale-$k$ units the non-renormalising coefficient $c_k=t_k^2(\dm_{\mathrm{fin}}+\dm_{K,k})$ of the \cite[(1.7)]{Ba3} vertex obeys $|c_k|\le C_Tt_k(1+\ln t_k^{-1})^5$ with $C_T$ free of $\eps_0$ (\cref{thm:a06e:T0}; asserted on \cite[p.~418]{Ba3} as $\dm_{K,k}t_k^2=O(t_k)$, $\dm_{\mathrm{fin}}=O(1)$, the first of which is \cref{thm:a06e:DM}), entering only through Proposition~1's order count (\cite[p.~421]{Ba3}: ``the vertices (1.7) give also a positive power of $L^k\eps$'') and in kind (b3) of \cref{lem:a06d:V7}; every bound of Steps~3(D$_1$)--3(D$_2$) stated with a factor $t_k$ from such a vertex holds with $t_k(1+\ln t_k^{-1})^5$ in its place, which its consumers absorb (\cref{lem:a06e:T1}). Geometry: \cite[(2.7)--(2.8)]{Ba2} and the bound $e^{-\delta c_*R(1+\log t^{-1})^\varrho}\le C_\kappa t^\kappa$ for every $\kappa$, uniformly in $t\in(0,1]$ ($\varrho>1$; \cref{lem:a06A:polylog}).

\begin{lemma}[gluing two families along contractions]\label{lem:a06d:V10}
Let $F_1,F_2$ be families obeying \cite[(1.33)]{Ba3} with multi-variable external fields on their legs, $|E(F_i,\Phi)|\le P_ie^{-\delta_0d_i}\norm{h\Phi}_{1,\alpha_0}$, and let $K_1,\dots,K_r$ be two-variable kernels joining a leg of $F_1$ (cube $\square_j$) to a leg of $F_2$ (cube $\square'_j$), with $\norm{h_{\square_j}K_jh_{\square'_j}}_{(1,\alpha_0)\otimes(1,\alpha_0)}\le\kappa_je^{-\delta'\dist(\square_j,\square'_j)}$ (the $C^{1,\alpha_0}$-norm in each variable uniformly in the other, with the mixed derivative and H\"older terms). Then $|E(F_1\circ_KF_2,\Phi_{\mathrm{rest}})|\le P_1P_2\prod_j\kappa_j\,e^{-\delta_0(d_1+d_2)}e^{-\delta'\sum_j\dist(\square_j,\square'_j)}\norm{h\Phi_{\mathrm{rest}}}$.
\end{lemma}

\begin{proof}
$E(F_1\circ_KF_2)=E(F_1,\Psi\otimes\Phi_1)$ with $\Psi(x_1,\dots,x_r)\defeq E(F_2,K_1(x_1,\cdot)\otimes\dots\otimes K_r(x_r,\cdot)\otimes\Phi_2)$. $\Psi$ is linear in each $K_j(x_j,\cdot)$; its sup, derivative and H\"older quotient in $x_j$ are $E(F_2,\cdot)$ applied to $K_j(x_j,\cdot)$, $\partial_{x_j}K_j(x_j,\cdot)$ and the H\"older quotients of $K_j$ in $x_j$, each bounded by $P_2e^{-\delta_0d_2}$ times the mixed norm. So $\norm{h\Psi}_{1,\alpha_0}\le P_2e^{-\delta_0d_2}\prod_j\kappa_je^{-\delta'\dist}$; apply \cite[(1.33)]{Ba3} to $F_1$. For $m\ge3$ families (a cumulant $\kappa_m$) induct, joining $F_m$ to the group already joined.
\end{proof}

Proposition~1's many-variable norm (\cite[p.~420]{Ba3}: \cite[(1.32)]{Ba3} ``extends in a natural way'') is $C^{1,\alpha_0}$ per variable; mixed derivatives enter only through the joins of \cref{lem:a06d:V10}, where the lifted kernels and their deep differences factorise with one derivative on each lifting factor; $\delta G$ lines never serve as joining kernels, and self-contractions inside one family are many-variable external fields of that family (I1c).

\begin{lemma}[localisation sums]\label{lem:a06d:V11}
Let $N\le\bar N$ be the number of indices and $\delta>0$. Then $\sum_{\Lambda:\square_{\iota_0}=\square_0}e^{-\delta d(\Lambda)}\le S_N(\delta)\defeq N^{N-2}(e^\delta\sigma(\delta))^{N-1}$, uniformly in the torus, with $\sigma(\delta)\defeq\sup_z\sum_\square e^{-\delta\dist(\square,z)}$. Consequently, with $S\defeq\max_{N\le\bar N}S_N(\delta)$: for a set $Z$ of blocks, $\sum_{\Lambda:\text{designated cube in }Z}Ct_ke^{-\delta d(\Lambda)}\le C\bar NSt_k|Z|_k$, and $\sum_{\Lambda:\text{anchor in the vertex domain}}Ct_ke^{-\delta\rho}e^{-\delta d(\Lambda)}\le C\bar NSt_ke^{-\delta\rho}N_k$.
\end{lemma}

\begin{proof}
$e^{-\delta d}\le\sum_{\text{trees}}e^{-\delta L_T}$; there are $N^{N-2}$ trees (Cayley); sum the cubes from the leaves inward, each edge contributing $\le e^\delta\sigma(\delta)$.
\end{proof}

\begin{lemma}[bounds per family and per localisation, from Proposition~1]\label{lem:a06d:V9}
There are $C_P<\infty$ and $\delta>0$, depending only on $(d,N,\norm q,a,L,M,\bar n,\alpha_0,R,C_g)$ and on the constants of (I1)--(I3), not on $\eps$, $k$, the torus or the data, such that for every localisation $\Lambda$: \textup{(a)} every grouped contour family ($\mathfrak B$-type, \cref{lem:a06d:K2}) on any datum obeys $|\cdot(\Lambda)|\le C_Pt_ke^{-\delta d(\Lambda)}$; \textup{(b)} a grouped $\Delta$-type family, or any grouped family with a designated $\delta G$ line or a designated $\delta$ in the external-field map, whose designated object sits at or joins cubes at depth $\ge\rho\ge R_0+2$ from the boundary where the two propagators differ, obeys $|\cdot(\Lambda)|\le C_Pt_ke^{-\delta\rho}e^{-\delta d(\Lambda)}$; \textup{(c)} every grouped step-$k$ cumulant family (a term of $\mathrm{out}_k$ on any datum) obeys $|\cdot(\Lambda)|\le C_Pt_ke^{-\delta d(\Lambda)}$, and with a designated $\delta C$ contraction at depth $\ge\rho$ it obeys $|\cdot(\Lambda)|\le C_Pt_ke^{-\delta\rho}e^{-\delta d(\Lambda)}$; \textup{(d)} the families of (a)--(c) with one designated vertex of kind (b1), (b2) or (b3) of \cref{lem:a06d:V7} obey the same bounds, and their $\Lambda$ contains a cube that meets the zone $Z$.
\end{lemma}

\begin{proof}
\emph{Scope.} Each family lies in the scope of (I1): (a) is (I1a) without external field; (b) is (I1b) with the designated $\delta G$ as a two-variable external field; (c) is (I1c), with \cref{lem:a06d:V10} when a cumulant joins several families; (d) is (I1d): for (b1) the designated leg carries the multiplier $(g'-g)$, in the class and supported in $Z$; for (b2) the scalar vertices are localised by unit cubes anyway \cite[p.~420]{Ba3}; for (b3) the vertex is a finite mass vertex with coefficient $\le Ct_k$ (I3), in Proposition~1's class inside the families of \cref{lem:a06d:V7} (its loop line is $\delta_k$, a lifted index-$k$ line, or not a pure full-line loop), never alone. (A lone $\mathrm{out}$-term is not bounded either: a figure-eight with one full and one index-$k$ line is $\sim\lambda t\eta^{-1}$; the pairs in $\mathrm{out}_k$ are inherited from the renormalised vacuum part of $F_k$.) \emph{Order.} Every vacuum family has order $o\defeq d_v/2+d_s\ge1$: all $A'$-legs are contracted, so $d_v$ is even, and if $d_v=0$ the family contains a \cite[(1.6)]{Ba3} vertex or a \cite[(1.7)]{Ba3} vertex of positive order; so $e(t_k)^{d_v}\lambda(t_k)^{d_s}=e^{d_v}\lambda^{d_s}t_k^o\le Ct_k$ for $t_k\le1$; a step-$k$ cumulant family is a join of families of order $\ge1$. \emph{External fields.} In (b) and (c), (I2a)--(I2c) bound the designated objects by $Ce^{-(\delta'/4)(\rho_\square+\rho_{\square'}+\dist(\square,\square'))}$ with $\rho_\square,\rho_{\square'}\ge\rho$, so $e^{-\delta'(\rho_\square+\rho_{\square'})/4}\le e^{-\delta'\rho/2}$ and $e^{-\delta'\dist(\square,\square')/4}$ supplies the tree link between the two cubes; the other contractions in (c) are lifted index-$k$ kernels with $\norm{h\cdot h'}\le Ce^{-\delta'\dist}$ (I2c). \emph{Assembly.} \cite[(1.33)]{Ba3} gives $O(1)Ct_ke^{-\delta_0d}\prod(\text{norms})$; the tree length over all cubes is at most the sum of the constituent tree lengths and the link lengths. Take $\delta\defeq\frac12\min(\delta_0,\delta'/4)$ and $C_P\defeq(\text{number of families of order}\le\bar n)\cdot O(1)\cdot C$.
\end{proof}

\begin{theorem}[$(\mathrm{VAC}^\flat)_k$, $0\le k\le K$]\label{thm:a06d:V12}
Let $\mathfrak D$ be a datum at scale $k$ with $\rho(\mathfrak D)\ge R_0+3$. Then for every localisation $\Lambda$, $|\mathcal R_k[\mathfrak D](\Lambda)|\le C_Vt_ke^{-\delta d(\Lambda)}$ with $C_V\defeq2N_{\mathrm{fam}}C_P$, $N_{\mathrm{fam}}$ the number of families of order $\le\bar n$; consequently $|\mathcal R_k[\mathfrak D]|\le C_V\bar NSt_k|\Omega_1|_k$. At $\mathfrak D=\mathfrak D_K$ this gives $|\mathcal R_K|\le C_V\bar NSt_K|T|_K=C_V\bar NSt_K^{-2}|T|\le C_V\bar NSL^2\eps_0^{-2}|T|$, since $t_K\in(\eps_0/L,\eps_0]$; at $\mathfrak D=\mathfrak D_k$ the hypothesis holds, $\rho(\mathfrak D_k)\ge r(t_k)-c_d\ge R-c_d\ge R_0+3$ (\cref{lem:a06d:V8}(c) at $L=5$), given $R\ge R_0+3+c_d$ (admissible: $R$ is chosen after $M$, \cite[(2.7)]{Ba2}).
\end{theorem}

\begin{proof}
$\mathcal R_k=\mathfrak B_k+\Delta_k$ (\cref{lem:a06d:V5}). $\mathfrak B_k$ is bounded by \cref{lem:a06d:V9}(a). $\Delta_k$ is bounded by \cref{lem:a06d:V9}(b) through its grouped families (\cref{lem:a06d:K2}): the designated $\delta_k=G_k(\Omega,0)-G_k(T,0)$ sits between vertices of $\Omega_1$ plus one block, at depth $\ge\rho(\mathfrak D)-1\ge R_0+2$, as (b) requires. Sum over localisations with \cref{lem:a06d:V11}.
\end{proof}

\begin{corollary}[{the vacuum clause of \cite[Prop.~2.1, (2.57)]{Ba2} at zero background}]\label{cor:a06d:V13}
Let $\mathfrak D'\defeq(\Omega^{(k)},B^k(\Lambda^{(k)}_7),g',c')$ with $g'$ admissible for $B^k(\Lambda^{(k)}_7)$ and $c'$ its \cite[(1.29)]{Ba3}. Then $\mathcal R_k[\mathfrak D_k]-\mathcal R_k[\mathfrak D']$ is a sum of families with a designated vertex in $Z_{\mathrm{col}}\defeq B^k(\mathrm{Col}_k)\cup\{g'<1\}$, and $|\mathcal R_k[\mathfrak D_k]-\mathcal R_k[\mathfrak D']|\le Ct_k|\mathrm{Col}_k|_k$ with $C\defeq C_P\bar NSC(\delta)(1+(2M+5)^3)$.
\end{corollary}

At $A^{(k)}=0$, $\phi=0$ the quartic term of \cite[(2.57)]{Ba2} vanishes and $\Efin$, $E_{K,k}$ are unchanged, so this is the $c$-number clause of \cite[(2.57)]{Ba2}, with $O(t_k)\le O(t_k^{\kappa_0})$ per collar block, under \cref{conv:a06d:LOC} (the mass counterterm enters as a vertex coefficient, and the replacement of propagators acts on lines only).

\begin{proof}
Apply \cref{lem:a06d:V7}(b): the vertex domains differ by $\Omega_1\triangle\Omega'_1=B^k(\mathrm{Col}_k)$; on $B^k(\Lambda^{(k)}_7)$, $g_k=1$ (\cref{lem:a06d:V8}(d)), so $\supp(g_k-g')\subset B^k(\mathrm{Col}_k)\cup\{g'<1\}$ and $c_k-c'$ is supported there up to tails $e^{-\delta(7r-M)}$. Apply \cref{lem:a06d:V9}(d), then \cref{lem:a06d:V8}(f) with $D_Z=M+1$, then \cref{lem:a06d:V11}; the coefficient tails are summed by $\sum_xe^{-\delta\dist(x,Z)}\le C(\delta)|Z|$.
\end{proof}
